\chapterbackmatter{Solutions to Supplementary Exercises}

\begin{enumerate}
\SupplementarySolution{1}{The Fourth Connected Moment}
Write \(m_B=\langle\prod_{i\in B}X_i\rangle\).  The partition relation
contains one four-element block, four partitions of type \(3+1\), three
of type \(2+2\), six of type \(2+1+1\), and one of type \(1+1+1+1\).
Solving recursively gives
\[
\begin{split}
 \langle X_1X_2X_3X_4\rangle_c={}&m_{1234}
 -\sum_{\rm 4\ triples}m_{ijk}m_l
 -\bigl(m_{12}m_{34}+m_{13}m_{24}+m_{14}m_{23}\bigr)\\
 &+2\sum_{\rm 6\ pairs}m_{ij}m_km_l
 -6m_1m_2m_3m_4.
\end{split}
\]
For centered variables the terms containing one-point moments vanish,
leaving the familiar subtraction of the three pairwise products.  The
coefficients \(1,-1,-1,2,-6\) are the partition-lattice M{\o}bius
coefficients.

\SupplementarySolution{2}{M{\o}bius Inversion on Set Partitions}
Insert the proposed inverse into the moment formula.  The coefficient of
a fixed product \(\prod_{B\in\rho}m(B)\) is a sum over partitions of the
\(|\rho|\) blocks into larger groups:
\[
 \sum_{\sigma\in\Pi(\rho)}
 (-1)^{|\sigma|-1}(|\sigma|-1)!.
\]
When \(|\rho|=1\), this is one.  When \(|\rho|>1\), group terms by the
block containing a chosen element, or equivalently use
\[
 \sum_{k=1}^{n}S(n,k)(-1)^{k-1}(k-1)!=\delta_{n1},
\]
where \(S(n,k)\) is a Stirling number of the second kind.  Thus every
product associated with a proper partition cancels, while the
single-block moment remains.  This proves
\[
 \kappa(A)=\sum_{\pi\in\Pi(A)}
 (-1)^{|\pi|-1}(|\pi|-1)!\prod_{B\in\pi}m(B).
\]

\SupplementarySolution{3}{Mixed Cumulants of Independent Clusters}
Introduce sources \(t_i\) for the \(X\)'s and \(s_j\) for the \(Y\)'s.
Independence gives
\[
 Z(t,s)=\left\langle
 e^{\sum_i t_iX_i+\sum_j s_jY_j}\right\rangle
 =Z_X(t)Z_Y(s).
\]
Therefore the connected generating function is
\[
 W(t,s)=\log Z(t,s)=W_X(t)+W_Y(s).
\]
Any mixed connected moment differentiates \(W\) at least once with
respect to a \(t_i\) and once with respect to an \(s_j\), and is therefore
zero.  Cluster decomposition is the spacetime version of this statement:
when separated experiments become independent, connected quantities
linking the two regions vanish, while full quantities factorize.

\SupplementarySolution{4}{Why the Logarithm Selects Connected Diagrams}
If a disconnected graph contains \(m_C\) indistinguishable copies of
component \(C\), permutations of these copies do not produce a new graph.
Its value is therefore
\[
 V_D=\prod_C\frac{V_C^{m_C}}{m_C!}.
\]
Summing independently over every nonnegative \(m_C\),
\[
 \sum_DV_D
 =\prod_C\sum_{m_C=0}^\infty\frac{V_C^{m_C}}{m_C!}
 =\prod_Ce^{V_C}
 =\exp\!\left(\sum_CV_C\right).
\]
Taking the logarithm removes all products and leaves exactly one copy of
each connected component value.

\SupplementarySolution{5}{Connected Derivatives of a Source Functional}
Use condensed indices and write
\[
 Z[J]=1+Z_1+\frac1{2!}Z_2+\frac1{3!}Z_3+\frac1{4!}Z_4+\cdots,
\]
where \(Z_n\) is homogeneous of degree \(n\) in \(J\).  Expanding
\(\log Z\) gives
\[
 W_1=Z_1,\qquad
 W_2=Z_2-Z_1^2,
\]
\[
 W_3=Z_3-3Z_2Z_1+2Z_1^3,
\]
with the products understood as sums over assignments of labeled source
indices.  At fourth order,
\[
 W_4=Z_4-4Z_3Z_1-3Z_2^2+12Z_2Z_1^2-6Z_1^4.
\]
Functional differentiation assigns the external labels to every
displayed factor and reproduces the partition formulas of S.4.1.
Consequently
\[
 \left.\frac{\delta^nW}{\delta J_1\cdots\delta J_n}\right|_{J=0}
\]
is the connected \(n\)-point function, while the same derivative of
\(Z\) is the full one.

\SupplementarySolution{6}{Volume Scaling of Connected Pieces}
In a periodic box,
\[
 \int_Vd^3x\,e^{i(\mathbf P_f-\mathbf P_i)\cdot\mathbf x}
 =V\,\delta_{\mathbf P_f,\mathbf P_i}.
\]
A connected component has one freely translatable center and therefore
one such factor.  Its remaining integrals involve relative coordinates
and do not acquire another overall power of \(V\).  A graph with \(C\)
disconnected components has \(C\) independently translatable centers,
so its leading volume dependence is \(V^C\).

Since the full generating object is the exponential of its connected
pieces, its logarithm contains only one connected component at a time and
is proportional to \(V\).  This is the linked-cluster origin of
extensivity.

\SupplementarySolution{7}{Cancellation of Vacuum Bubbles}
Let \(B\) be the sum of connected vacuum bubbles.  Diagram
exponentiation gives
\[
 \bra0S\ket0=e^B.
\]
In the numerator with external fields, every diagram uniquely splits
into vacuum components and a remainder in which each component touches
at least one external point.  Hence
\[
 \bra0T\{\phi_1\cdots\phi_nS\}\ket0
 =e^B\,\mathcal N_n.
\]
The Gell-Mann--Low ratio is therefore \(\mathcal N_n\): the factor \(e^B\)
cancels exactly.

A vacuum bubble touches no external point.  By contrast, a diagram may
have two separate components, each attached to some external points.
Such a diagram survives vacuum normalization but is removed when one
passes from the full correlator to its connected part.

\SupplementarySolution{8}{Connected Is Not One-Particle-Irreducible}
A graph is connected if every pair of vertices is joined by some path.  It
is one-particle irreducible if it remains connected after any one
internal line is cut.  For example, take two interaction subgraphs joined
by a single internal propagator, with one external leg attached at each
end.  The complete graph is connected, but cutting the joining propagator
separates it into two pieces.

Connected correlation functions sum all connected graphs, including this
one-particle-reducible example.  One-particle-irreducible graphs form a
smaller set used to build effective vertices; chains of them reconstruct
general connected graphs.

\SupplementarySolution{9}{One Delta Function per Connected Component}
Choose a spanning forest, one spanning tree in each connected component.
It contains \(V-C\) edges.  Each of the remaining
\[
 L=I-(V-C)
\]
internal edges creates one independent cycle.  Rearrangement gives
\[
 V-I+L=C.
\]
Every vertex initially supplies one momentum delta function.  Of the
\(I\) internal momenta, \(I-L\) are fixed by these constraints; the
\(L\) loop momenta remain to be integrated.  The number of unused delta
functions is therefore
\[
 V-(I-L)=V-I+L=C.
\]
They impose conservation of the total momentum entering each connected
component.  A connected graph has \(C=1\) and hence only the overall
conservation delta function.

\SupplementarySolution{10}{Locality Produces a Single Vertex Delta Function}
Each field mode contributes a plane wave.  A typical operator term at
fixed time is proportional to
\[
 \int d^3x\,
 e^{i(\mathbf p'_1+\cdots+\mathbf p'_N
 -\mathbf p_1-\cdots-\mathbf p_M)\cdot\mathbf x}.
\]
The common spatial argument of all four fields makes the integral
\[
 (2\pi)^3\delta^3\!\left(
 \sum_{r=1}^N\mathbf p'_r-\sum_{s=1}^M\mathbf p_s\right).
\]
There is only one spatial integration variable, so there can be only one
independent momentum delta function.  Its coefficient contains mode
normalizations and operator products but no further delta support.

\SupplementarySolution{11}{A Bilocal Kernel and Extra Singular Support}
Fourier expand
\[
 K(\mathbf r)=\int\frac{d^3k}{(2\pi)^3}
 \widetilde K(\mathbf k)e^{i\mathbf k\cdot\mathbf r}.
\]
The \(x\) and \(y\) integrals give a total-momentum delta function and a
factor \(\widetilde K(\mathbf Q)\), where \(\mathbf Q\) is the momentum
transferred between the two pairs:
\[
 h\propto(2\pi)^3\delta^3(\mathbf P_f-\mathbf P_i)
 \widetilde K(\mathbf Q).
\]
If \(\widetilde K\) is an ordinary function, this has only the required
overall delta distribution.  For a local kernel
\(K(\mathbf r)\propto\delta^3(\mathbf r)\),
\(\widetilde K\) is constant.

If \(K(\mathbf r)=K_0\), however,
\(\widetilde K(\mathbf Q)=(2\pi)^3K_0\delta^3(\mathbf Q)\).  The vertex
then separately conserves the momentum of each pair and contains two
delta functions.  This infinite-range interaction correlates separated
clusters and violates the cluster condition.

\SupplementarySolution{12}{Four-Particle Cluster Factorization}
When only the clusters \(12\) and \(34\) remain mutually connected, the
surviving terms are
\[
\begin{split}
 S_{1'2'3'4',1234}\longrightarrow{}&
 S^C_{1'2',12}S^C_{3'4',34}\\
 &+I_{1'2',12}S^C_{3'4',34}
 +S^C_{1'2',12}I_{3'4',34}\\
 &+I_{1'2',12}I_{3'4',34},
\end{split}
\]
where \(I_{1'2',12}=\delta_{1'1}\delta_{2'2}\) for distinguishable
particles.  Terms connected across the two distant laboratories vanish.
Factoring the four surviving terms gives
\[
 (S^C_{1'2',12}+I_{1'2',12})
 (S^C_{3'4',34}+I_{3'4',34})
 =S_{1'2',12}S_{3'4',34},
\]
which is the cluster-decomposition statement.

\SupplementarySolution{13}{Wick's Theorem for Three Fields}
Write \(\phi_i=\phi_i^{(+)}+\phi_i^{(-)}\), with the positive-frequency
part annihilating the vacuum and the negative-frequency part creating
quanta.  Moving every annihilation part to the right produces one
commutator for each interchange with a creation part.  Since three fields
cannot be completely paired, the possibilities are no contraction or one
contraction:
\[
\begin{split}
 T\{\phi_1\phi_2\phi_3\}={}&:\phi_1\phi_2\phi_3:
 +D_F(x_1-x_2)\phi_3\\
 &+D_F(x_1-x_3)\phi_2
 +D_F(x_2-x_3)\phi_1.
\end{split}
\]
The propagators are c-numbers, so their position relative to the
uncontracted field is immaterial.  The result is symmetric in all three
points, as required by time ordering for a scalar field.

\SupplementarySolution{14}{The First Vacuum Graphs in Cubic Theory}
The zeroth-order term is one, and the first-order term vanishes because
it contains three free fields.  With the vertex factor denoted \(i\lambda\),
the second-order term is
\[
 \frac{(i\lambda)^2}{2(3!)^2}
 \int d^4x\,d^4y\,
 \langle0|T\{\phi(x)^3\phi(y)^3\}|0\rangle.
\]
There are \(3!=6\) contractions joining every field at \(x\) to one at
\(y\), and \(3\cdot3=9\) contractions with one joining line and one
tadpole at each vertex.  Thus
\[
 \bra0S\ket0
 =1+(i\lambda)^2\int d^4x\,d^4y
 \left[\frac{D_F(x-y)^3}{12}
 +\frac{D_F(x-y)D_F(0)^2}{8}\right]+O(\lambda^4).
\]
The two connected vacuum graphs have symmetry factors \(12\) and \(8\).

\SupplementarySolution{15}{Odd Orders and Vacuum Pairings}
At order \(n\), cubic theory contains \(3n\) free fields.  A vacuum
expectation value is a sum over complete pairings and hence vanishes when
\(3n\) is odd.  Since \(3n\) has the parity of \(n\), every odd
perturbative order vanishes.

At fourth order, a vacuum graph is either connected across all four
vertices or is a union of two connected two-vertex graphs.  If \(C_2\)
denotes the complete connected second-order contribution and \(C_4\) the
sum of connected fourth-order graphs, the order-four result has the
universal form
\[
 C_4+\frac12C_2^2.
\]
The square includes all pairings of the two possible second-order
topologies from S.4.14; the \(1/2!\) removes interchange of the two
disconnected components.

\SupplementarySolution{16}{Normalized Two-Point Functions}
Let \(G_0(x,y)\) be the free two-point function and \(B_2\) the complete
second-order vacuum contribution.  The numerator can be organized as
\[
 \mathcal N=G_0+\lambda^2
 \bigl(G_{2,\rm attached}+G_0B_2\bigr)+O(\lambda^4),
\]
where \(G_{2,\rm attached}\) contains precisely the terms in which every
interaction vertex connects to an external point.  The denominator is
\[
 \mathcal D=1+\lambda^2B_2+O(\lambda^4).
\]
Using \(\mathcal D^{-1}=1-\lambda^2B_2+\cdots\),
\[
 \frac{\mathcal N}{\mathcal D}
 =G_0+\lambda^2G_{2,\rm attached}+O(\lambda^4).
\]
The two copies of \(G_0B_2\) cancel algebraically.  The same factorization
and cancellation holds for any number of vacuum components and to every
order.

\SupplementarySolution{17}{External-Connected and Fully Connected}
Condition (i) excludes vacuum bubbles: each interaction vertex must lie
in a component that touches some external point.  It does not require all
external points to share one component.  For example, a four-point
diagram may be the product of two dressed propagators, one joining
\((x_1,x_2)\) and one joining \((x_3,x_4)\).  Every vertex touches an
external point, but the graph has two components.

Vacuum normalization retains this product.  The connected four-point
function subtracts
\[
 G(x_1,x_2)G(x_3,x_4)+G(x_1,x_3)G(x_2,x_4)
 +G(x_1,x_4)G(x_2,x_3),
\]
along with one-point terms when present.  The two-component example is
therefore removed, leaving only graphs in which all four external points
belong to one component.

\SupplementarySolution{18}{A Symmetry Factor from Wick Contractions}
The second Dyson term supplies \(1/[2!(3!)^2]=1/72\).  For the theta graph,
all three fields at one vertex contract with the other vertex, giving
\(3!=6\) pairings.  Its coefficient is \(6/72=1/12\), so its symmetry
factor is \(12\).

For the tadpole graph, choose the field on the first vertex that joins the
second in three ways and its partner there in three ways.  The remaining
two fields at each vertex pair uniquely, for \(3\cdot3=9\) contractions.
Its coefficient is \(9/72=1/8\), so its symmetry factor is \(8\).  These
counts agree with the automorphism groups of the two graphs.

\SupplementarySolution{19}{A Normalized Coherent Eigenstate}
The normalized state is
\[
 \ket z=e^{-|z|^2/2}e^{za^\dagger}\ket0
 =e^{-|z|^2/2}\sum_{n=0}^\infty
 \frac{z^n}{\sqrt{n!}}\ket n.
\]
Since \([a,za^\dagger]=z\),
\(ae^{za^\dagger}=e^{za^\dagger}(a+z)\), so
\(a\ket z=z\ket z\).  Its overlap is
\[
 \braket{z_2}{z_1}
 =\exp\!\left(-\frac{|z_2|^2}{2}-\frac{|z_1|^2}{2}
 +z_2^*z_1\right).
\]
Setting \(z_2=z_1\) verifies unit normalization; distinct coherent states
are generally not orthogonal.

\SupplementarySolution{20}{Resolution of the Identity by Coherent States}
Insert the number-state expansions:
\[
 \int\frac{d^2z}{\pi}\ket z\bra z
 =\sum_{m,n}\frac{\ket n\bra m}{\sqrt{n!m!}}
 \int\frac{d^2z}{\pi}e^{-|z|^2}z^n(z^*)^m.
\]
The angular integral vanishes unless \(m=n\), and the radial integral is
\[
 2\int_0^\infty r\,dr\,e^{-r^2}r^{2n}=n!.
\]
Thus the operator is \(\sum_n\ket n\bra n=\mathbf1\).  Orthogonality is
not required for a resolution of the identity: the continuous family is
overcomplete, and its integration measure compensates for the nonzero
overlaps.

\SupplementarySolution{21}{Means and Minimal Uncertainty}
Using \(a\ket z=z\ket z\),
\[
 \langle q\rangle=\frac{z+z^*}{\sqrt2}
 =\sqrt2\,\operatorname{Re}z,\qquad
 \langle p\rangle=\frac{z-z^*}{i\sqrt2}
 =\sqrt2\,\operatorname{Im}z.
\]
Normal ordering \(q^2\) and \(p^2\) and using
\([a,a^\dagger]=1\) gives
\[
 \langle q^2\rangle=\langle q\rangle^2+\frac12,\qquad
 \langle p^2\rangle=\langle p\rangle^2+\frac12.
\]
Hence
\[
 \Delta q=\Delta p=\frac1{\sqrt2},\qquad
 \Delta q\,\Delta p=\frac12.
\]
The displacement changes the means but leaves the vacuum fluctuations
unchanged.

\SupplementarySolution{22}{Multimode Number Statistics}
Let
\(\mu=\int dq\,|\lambda(q)|^2\).  Expanding the coherent state by total
particle number gives a Poisson distribution, so its moment-generating
function is
\[
 \langle e^{tN}\rangle
 =e^{-\mu}\sum_{n=0}^\infty
 \frac{\mu^n e^{tn}}{n!}
 =\exp\!\bigl[\mu(e^t-1)\bigr].
\]
Differentiation at \(t=0\) yields
\[
 \langle N\rangle=\mu,\qquad
 \operatorname{Var}N=\mu,
\qquad
 P(N=n)=e^{-\mu}\frac{\mu^n}{n!}.
\]
Disjoint momentum bins are independent Poisson variables, another
manifestation of the exponential connected-state structure.

\SupplementarySolution{23}{The Displacement Operator}
Put \(A=za^\dagger\) and \(B=-z^*a\).  Their commutator
\([A,B]=|z|^2\) is a c-number, so the
Baker--Campbell--Hausdorff formula gives
\[
 D(z)=e^{A+B}=e^{-|z|^2/2}e^{za^\dagger}e^{-z^*a}.
\]
Since the exponent \(za^\dagger-z^*a\) is anti-Hermitian,
\(D(z)^\dagger=D(-z)=D(z)^{-1}\).  Acting on the vacuum,
\[
 D(z)\ket0=e^{-|z|^2/2}e^{za^\dagger}\ket0=\ket z.
\]
Conjugation also gives \(D(z)^\dagger aD(z)=a+z\), directly confirming
the coherent-state eigenvalue.

\SupplementarySolution{24}{In- and Out-Vacua in a Classical Source}
The relation in the question implies
\[
 a_{\rm in}(\mathbf p)
 =a_{\rm out}(\mathbf p)-\alpha(\mathbf p).
\]
Therefore \(a_{\rm in}\InKet{0}=0\) is equivalent to
\[
 a_{\rm out}(\mathbf p)\InKet{0}
 =\alpha(\mathbf p)\InKet{0}.
\]
Thus, with \(\mu=\int d^3p\,|\alpha(\mathbf p)|^2\),
\[
 \InKet{0}
 =e^{-\mu/2}
 \exp\!\left[\int d^3p\,\alpha(\mathbf p)
 a_{\rm out}^\dagger(\mathbf p)\right]\OutKet{0}.
\]
The vacuum persistence amplitude is
\(\OutBra{0}\InKet{0}=e^{-\mu/2}\), so the probability of producing no
out-particles is \(P_0=e^{-\mu}\).

\SupplementarySolution{25}{An Impulse Creates a Coherent Field}
Choose
\[
 \phi(\mathbf x)=\int\frac{d^3p}{(2\pi)^{3/2}}
 \frac{a(\mathbf p)e^{i\mathbf p\cdot\mathbf x}
 +a^\dagger(\mathbf p)e^{-i\mathbf p\cdot\mathbf x}}
 {\sqrt{2\omega_{\mathbf p}}},
\quad \omega_{\mathbf p}=\sqrt{\mathbf p^2+m^2}.
\]
Integrating the equations across \(t=0\) gives
\[
 \phi(0^+)=\phi(0^-),\qquad
 \pi(0^+)=\pi(0^-)+j.
\]
Inverting the mode expansion,
\[
 a(\mathbf p,0^+)=a(\mathbf p,0^-)
 +\frac{i\,\widetilde j(\mathbf p)}
 {\sqrt{2\omega_{\mathbf p}}}.
\]
If the state was annihilated by \(a(\mathbf p,0^-)\), then it is an
eigenstate of the post-impulse operator with eigenvalue
\[
 \alpha(\mathbf p)=
 \frac{i\,\widetilde j(\mathbf p)}
 {\sqrt{2\omega_{\mathbf p}}},
\]
up to the displayed Fourier normalization.  It is therefore a multimode
coherent state.

\SupplementarySolution{26}{Classical Expectation Values from a Coherent State}
For \(a(\mathbf p)\ket\lambda=\lambda(\mathbf p)\ket\lambda\),
\[
 \langle\phi(x)\rangle_\lambda
 =\int\frac{d^3p}{(2\pi)^{3/2}\sqrt{2\omega_{\mathbf p}}}
 \left[\lambda(\mathbf p)e^{-i\omega t+i\mathbf p\cdot\mathbf x}
 +\lambda(\mathbf p)^*e^{i\omega t-i\mathbf p\cdot\mathbf x}\right].
\]
Every mode is on shell, so
\((\Box-m^2)\langle\phi\rangle_\lambda=0\) in the local mostly-plus
convention.

Write \(a=\lambda+\delta a\).  The fluctuation operators \(\delta a\)
annihilate the displaced vacuum and have the same contractions as the
ordinary vacuum.  Hence
\[
 \langle\phi(x)\phi(y)\rangle_\lambda
 -\langle\phi(x)\rangle_\lambda\langle\phi(y)\rangle_\lambda
 =\langle0|\phi(x)\phi(y)|0\rangle.
\]
The state carries a classical mean field but unchanged connected
Gaussian fluctuations.

\SupplementarySolution{27}{The Momentum Spectrum of Source Radiation}
From S.4.25,
\[
 \langle a^\dagger(\mathbf p)a(\mathbf q)\rangle
 =\alpha(\mathbf p)^*\alpha(\mathbf q),
\qquad
 |\alpha(\mathbf p)|^2
 =\frac{|\widetilde j(\mathbf p)|^2}
 {2\omega_{\mathbf p}}.
\]
Thus the total mean number is
\[
 \mu=\int d^3p\,\frac{|\widetilde j(\mathbf p)|^2}
 {2\omega_{\mathbf p}},
\]
with the matching Fourier-measure factors for the chosen convention.
Coherent-state statistics gives \(P_0=e^{-\mu}\).  The state is a vector
in Fock space precisely when \(\mu<\infty\); otherwise its formal
displacement has infinite norm and belongs to an inequivalent
representation.

\SupplementarySolution{28}{Correlation Is Not Signaling}
For spacelike separation in the mostly-plus metric, put
\(r=\sqrt{(x-y)^2}\).  The free massive Wightman function is
proportional to
\[
 \frac{m}{4\pi^2r}K_1(mr),
\]
which is nonzero but behaves as \(e^{-mr}\) at large \(r\).  It measures
correlations already present in the vacuum.

The commutator is the difference of the two orderings,
\[
 [\phi(x),\phi(y)]=D^+(x-y)-D^+(y-x).
\]
Lorentz invariance allows a frame in which spacelike-separated points
are simultaneous; the two terms are then equal, so the commutator
vanishes.  Linear response is governed by the retarded function
\(i\theta(x^0-y^0)[\phi(x),\phi(y)]\), not by the Wightman correlation.
It therefore has causal support even though the vacuum is correlated.

\SupplementarySolution{29}{Separated Sources and Additivity of \(W\)}
The source expansion groups into connected components.  If connected
propagation between regions \(A\) and \(B\) is negligible, no connected
component contains insertions from both sources.  Hence the connected
functional splits:
\[
 W[J_A+J_B]\longrightarrow W[J_A]+W[J_B].
\]
Exponentiating gives
\[
 Z[J_A+J_B]\longrightarrow
 e^{W[J_A]}e^{W[J_B]}=Z[J_A]Z[J_B].
\]
Functional derivatives with respect to sources in both regions therefore
have vanishing connected part, while full transition amplitudes and
probabilities factorize.  This is precisely the statement that the two
distant experiments become independent.

\SupplementarySolution{30}{Cluster Decay and Fourier Singularities}
Remove the overall delta function and let \(\mathbf q\) denote momentum
transferred between the two proposed clusters.  Translating one cluster
by \(\mathbf R\) multiplies the wave-packet integral by
\(e^{i\mathbf q\cdot\mathbf R}\):
\[
 A_C(\mathbf R)=\int d^3q\,
 f(\mathbf q,\ldots)e^{i\mathbf q\cdot\mathbf R}.
\]
If the reduced connected kernel is locally integrable and the wave
packets make it \(L^1\), the Riemann--Lebesgue lemma gives
\(A_C(\mathbf R)\to0\) as \(|\mathbf R|\to\infty\).

An extra factor \(\delta^3(\mathbf q)\) instead gives
\[
 A_C(\mathbf R)\supset f(\mathbf0,\ldots),
\]
independent of \(\mathbf R\).  It separately conserves the momentum of a
proper cluster and prevents connected decay.  This is why one overall
momentum delta function is compatible with cluster decomposition, while
additional partial-momentum delta functions are not.
\end{enumerate}
